\clearpage
\fancyhead[L]{\small CERES.}
\hypertarget{native-317-lower}{}
\begin{center}
\textbf{CERES.}
\par\vspace{4pt}\hrule\vspace{4pt}
\fontsize{10}{12}\selectfont Göttingische gelehrte Anzeigen. Stück 87. S. 857..860. 1809. Juni 3.
\par\vspace{4pt}\hrule\vspace{4pt}
\end{center}

\par\indent Im 40. Stück unserer vorjährigen Anzeigen [S. 299 dieses Bandes] haben wir die zwölften Elemente der \emph{Ceres} mitgeteilt, die Herr Professor \textsc{Gauss} auf die Beobachtungen von 1807 gestützt hatte. Nun teilen wir die Ergebnisse der soeben abgeschlossenen Rechnungen aus den Beobachtungen von 1808 mit. Weil der Planet tief im Steinbock stand, war er in unseren nördlichen Gegenden nur sehr schwer zu beobachten. An der hiesigen Sternwarte war nur eine einzige Meridianbeobachtung möglich, die zudem nicht als besonders gut gelten kann.


\begin{center}\fontsize{10}{12}\selectfont\setlength{\tabcolsep}{5pt}\renewcommand{\arraystretch}{1.15}\begin{tabular}{lr} 1808 mittlere Zeit Juli 25. & $13^{\text{u}}\;14^{\text{m}}\;29^{\text{s}}$.\\ Rektaszension & $322^\circ\;8'\;57''6$.\\ Südliche Deklination & $27^\circ\;52'\;12''5$\\
\end{tabular}\end{center}


\SourceLeaf{318}\par\indent Weitere Beobachtungen sind bisher nur von Herrn von \textsc{Lindenau} an der Seeberger Sternwarte und von Herrn \textsc{Santini} in Padua bekannt geworden. Mithilfe dieser Beobachtungen hat Herr Professor \textsc{Gauss} die Opposition von 1808, den Fehler der XII. Elemente und die Korrekturen bestimmt, die an diesen Elementen anzubringen sind, um sie mit allen sechs bisher beobachteten Oppositionen möglichst genau in Übereinstimmung zu bringen. Da der Briefwechsel zurzeit so sehr erschwert ist, hoffen wir, dass es den Astronomen willkommen sein wird, diese Ergebnisse in unseren Blättern zu finden.

\par\indent Die Opposition wurde nach möglichst sorgfältiger Auswertung der Beobachtungen bestimmt. Das Ergebnis lautet wie folgt.

\begin{center} 1808 Aug. 5. \quad $11^{\text{u}}\ 30^{\text{m}}\ 13^{\text{s}}$ mittlere Zeit in Göttingen

wahre Länge der Ceres .\;.\;.\;\ $313^\circ\ 13'\ 40''2$

wahre geozentrische Breite .\;.\;\ $12\ 43\ 44{,}8$ südlich
\end{center}

\par\indent Die XII. Elemente weichen diesmal nur um eine Minute in der Länge und um ein paar Sekunden in der Breite ab. Das ist ein ganz unbedeutender Unterschied, wenn man bedenkt, dass die berücksichtigten Störungen noch nicht vollständig sind und daher eine vollkommene Übereinstimmung aller Beobachtungen aus fast acht Jahren ohnehin als unmöglich gelten muss. Herr Professor \textsc{Gauss} hat für die künftige vollständigere Berechnung der Störungen bereits alles Nötige entworfen, kann aber freilich noch nicht sagen, wann er diese ebenso heikle wie umfangreiche Arbeit tatsächlich ausführen kann. Inzwischen hielt er es für interessant zu untersuchen, wie genau sich alle sechs beobachteten Oppositionen noch mit seinen 1802 berechneten Störungstafeln in Einklang bringen lassen. Das Ergebnis zeigt, dass die Unterschiede auch diesmal noch so gering sind, dass man ihretwegen noch eine Zeit lang auf die vollständigeren Störungen verzichten kann. Man darf dabei aber nicht vergessen, dass die Berechnung einer einzigen Planetenposition gewiss viermal so viel Zeit und Arbeit kosten wird wie jetzt, sobald diese Störungen vollständig berücksichtigt werden müssen. Es folgen die neuen Elemente der Ceres, die nur äußerst wenig von den XII. abweichen. Die Epochen gelten für den Meridian von Göttingen.

\Needspace{23\baselineskip}
\begin{center} XIII. \emph{Elemente der Ceres.}
\end{center}


\begin{center}\fontsize{10}{12}\selectfont\setlength{\tabcolsep}{5pt}\renewcommand{\arraystretch}{1.15}\begin{tabular}{crr}  & Mittlere Länge. & Perihel.\\\hline 1801 & $77^\circ\;18'\;36''5$ & $146^\circ\;26'\;0''1$\\ 1802 & $155\;28\;23.4$ & $\phantom{146}\;28\;1.4$\\ 1803 & $233\;38\;10.3$ & $\phantom{146}\;30\;2.6$\\ 1804 & $312\;0\;48.1$ & $\phantom{146}\;32\;4.2$\\ 1805 & $30\;10\;35.0$ & $\phantom{146}\;34\;5.4$\\ 1806 & $108\;20\;21.8$ & $\phantom{146}\;36\;6.6$\\ 1807 & $186\;30\;8.7$ & $\phantom{146}\;38\;7.9$\\ 1808 & $264\;52\;46.5$ & $\phantom{146}\;40\;9.5$\\ 1809 & $343\;2\;33.4$ & $\phantom{146}\;42\;10.7$\\
\end{tabular}\end{center}

\begin{center}\fontsize{9}{11}\selectfont\setlength{\tabcolsep}{5pt}\renewcommand{\arraystretch}{1.15}\begin{tabular}{lr} Tägliche mittlere tropische Bewegung .\; & $770''9230$\\ Exzentrizität 1806 .\;.\;.\;.\;.\;.\;.\;.\; & $0.0785028$\\ Jährliche Abnahme .\;.\;.\;.\;.\;.\;.\;.\; & $0.00000583$\\ Logarithmus der großen Halbachse .\;.\; & $0.4420486$\\ Aufsteigender Knoten 1806 .\;.\;.\;.\; & $80^\circ\;53'\;41''3$\\ Jährliche Bewegung .\;.\;.\;.\;.\;.\;.\;.\; & $+1{,}48$\\ Neigung der Bahn 1806 .\;.\;,\;.\;.\; & $10^\circ\;37'\;31''2$\\ Jährliche Abnahme .\;.\;.\;.\;.\;.\;.\;.\;.\;.\; & $0{,}44$\\
\end{tabular}\end{center}


\Needspace{13\baselineskip}\par\indent Wie gut diese neuen Elemente mit allen sechs bisher beobachteten Oppositionen übereinstimmen, zeigen die folgenden Unterschiede.

\SourceLeaf{319}\begin{center} Unterschied:
\end{center}


\begin{center}\fontsize{10}{12}\selectfont\setlength{\tabcolsep}{5pt}\renewcommand{\arraystretch}{1.15}\begin{tabular}{crr}  & \shortstack{heliozentrische Länge\\in der Bahn} & \shortstack{heliozentrische\\Breite}\\\hline 1802 & $-0''5$ & $-1''9$\\ 1803 & $+22.1$ & $-5.3$\\ 1804 & $-34.9$ & $-1.2$\\ 1806 & $+26.7$ & $-3.6$\\ 1807 & $-26.9$ & $-11.7$\\ 1808 & $+20.8$ & $-6.0$\\
\end{tabular}\end{center}


\Needspace{17\baselineskip}\par\noindent Dabei ist zu bemerken, dass die Störungen der Länge und des Radiusvektors für diese Rechnung den Störungstafeln entnommen sind, die Herr Professor \textsc{Gauss} im Märzheft der Monatl. Corresp. [März 1803, Seite 235 und folgende Seiten dieses Bandes] veröffentlicht hat. Die Breitenstörungen beruhen dagegen auf einer späteren, vollständigeren Entwicklung, deren Ergebnisse bisher noch nicht gedruckt sind. Die geringen Unterschiede in der Breite darf man als Bestätigung dafür ansehen, dass bei jener Bestimmung keine erheblichen Glieder übersehen wurden. [Der Inhalt dieser Mitteilung ist in der Monatlichen Correspondenz, Band XIX, Seite 504\,..\,512, Mai 1809, wiedergegeben. Dort sind auch die von \textsc{Gauss} berechneten Ephemeriden abgedruckt, für Ceres vom 18. Juli 1809 bis zum 16. April 1810 und für Juno vom 4. September 1809 bis zum 10. Mai 1810. Sie geben für Mitternacht in Göttingen in Abständen von vier Tagen die geozentrische Rektaszension und die geozentrische Deklination auf die Minute genau und den Abstand von der Erde auf vier Dezimalstellen genau an.]

